\documentclass[12pt, a4paper]{article}

\usepackage[utf8]{inputenc}
\usepackage[T1]{fontenc}
\usepackage[serbian]{babel}
\usepackage[left=2.5cm, right=2.5cm, top=2.5cm, bottom=2.5cm]{geometry}
\usepackage{graphicx}
\usepackage{booktabs}
\usepackage{amsmath}
\usepackage{amssymb}
\usepackage{hyperref}
\usepackage{float}
\usepackage{caption}
\usepackage{subcaption}
\usepackage{array}
\usepackage{xcolor}
\usepackage{enumitem}
% microtype disabled: font expansion not available in this environment
\usepackage{parskip}

\hypersetup{
    colorlinks=true,
    linkcolor=black,
    urlcolor=blue,
    citecolor=black
}

\setlength{\parskip}{6pt}
\setlength{\parindent}{0pt}

\title{
    \textbf{Selekcija algoritma pretrage stabla igre}\\[0.4em]
    \large Istraživanje podataka 2: Seminarski rad\\[0.3em]
    \normalsize Primena klasifikacije na problem selekcije algoritma
}
\author{Anđela Spasić}
\date{Matematički fakultet, Univerzitet u Beogradu\\2025/2026}

\begin{document}

\maketitle
\newpage
\tableofcontents
\newpage

%-------------------------------------------------------------------
\section{Tekst zadatka}
%-------------------------------------------------------------------

Zadatak je primeniti metode istraživanja podataka na sopstveni sintetički skup stabala igre,
sa ciljem da se predvidi koji od dva algoritma pretrage, Alpha-Beta (AB) ili Monte Carlo Tree
Search (MCTS), daje bolji rezultat na datom stablu, pod ograničenim budžetom pretrage
(brojem čvorova koje algoritam sme da poseti).

Problem se formalizuje kao instanca \textit{problema selekcije algoritma} (Rice, 1976). Kao mera selekcije, u ovom radu koriste se dva glavna cilja klasifikacije
nad istim skupom:

\begin{itemize}[noitemsep]
    \item \texttt{which\_algo}: da li je algoritam $a$ pronašao globalno optimalan
    prvi potez na stablu $x$, prema iscrpnom minimax rešavaču;
    \item \texttt{which\_better}: čiji je od dva izabrana poteza (AB-ovog ili
    MCTS-ovog) vrednosno bolji, nezavisno od toga da li je bilo koji od njih zaista optimalan.
\end{itemize}

Cilj je naučiti model da predviđa
koji algoritam će dati bolji rezultat, pre nego što se ijedan od njih pokrene.

Optimalni potez i tačne vrednosti poteza definišu se iscrpnim minimax pretraživanjem koje
poseti svaki čvor stabla. Ovakvo pretraživanje je izvodljivo samo na sintetičkim stablima
kontrolisane veličine, pa se stabla generišu programski, upravo da bi iscrpni minimax mogao
da posluži za obeležavanje, tako da AB i MCTS imaju sa čime da se porede.

Seminarski rad obuhvata:
\begin{itemize}[noitemsep]
    \item generisani sintetički skup podataka (16\,978 stabala, 91 atribut, od čega 81
    prediktorski i 10 ciljnih/izvedenih)
    \item pretprocesiranje podataka
    \item klasifikaciju sa sedam algoritama i poređenje na tri skupa atributa (svi, strukturni,
    redukovani)
    \item vizualizaciju podataka u 2D prostoru (PCA, t-SNE)
    \item modele za više ciljnih promenljivih (\texttt{which\_better}, \texttt{which\_algo},
    \texttt{ab\_correct}, \texttt{mcts\_correct}, \texttt{margin\_category},
    \texttt{tree\_difficulty})
    \item tekstualni opis i tumačenje rezultata
\end{itemize}

%-------------------------------------------------------------------
\section{Opis skupa podataka}
%-------------------------------------------------------------------

\subsection{Generisanje stabala}

Svako stablo se gradi rekurzivno počev od korena. Vrednost svakog čvora živi u
neograničenom prostoru i predstavlja slučajni hod: dete nasleđuje vrednost roditelja uz mali Gausov pomak,
\[
\ell_{\text{dete}} = \ell_{\text{roditelj}} + \mathcal{N}(0, \sigma^2), \qquad
\sigma = \frac{\mathit{spread}}{\sqrt{\mathit{max\_depth}}},
\]
gde \textit{spread} kontroliše koliko brzo se vrednosti grana razilaze od korena ka listovima.
Nakon toga se vrednost projektuje u interval $[-1, 1]$ preko $\tanh$: list
dobija pravu vrednost $v = \tanh(\ell)$, a svaki unutrašnji čvor dobija \textit{heurističku}
vrednost $\tanh(\ell + \mathcal{N}(0, (\mathit{heur\_noise} \cdot r)^2))$, gde je $r$ udeo
preostale dubine do lista. Heuristika je namerno nepouzdanija bliže korenu, gde AB najčešće
mora da je koristi jer budžet ne dopire do pravih listova.

\textbf{Tipovi čvorova.} Koren uvek ima bar dva deteta (da postoji stvaran izbor poteza).
Svaki unutrašnji čvor je jedan od:
\begin{itemize}[noitemsep]
    \item \textsc{Max}: vrednost = $\max$ dece (naizmenično sa \textsc{Min} po parnosti dubine)
    \item \textsc{Min}: vrednost = $\min$ dece
    \item \textsc{Chance}: slučajni događaj, vrednost = $\text{mean}$ dece; javlja se na
    bilo kojoj dubini $> 0$ sa verovatnoćom \textit{chance\_density}, koja je za 15\% stabala
    tačno $0.0$ (potpuno deterministička stabla), a inače uniformno izvučena iz $[0, 0.9]$
    \item \textsc{Leaf}: terminalni čvor
\end{itemize}
Broj dece po čvoru je Gausov uzorak oko faktora grananja koji je jedan od parametara generisanja, a dodatno određeni broj čvorova (van korena)
postaje list pre dostizanja maksimalne dubine, čime stabla postaju nepravilna umesto savršeno
balansirana. Maksimalna dubina se bira tako da \textit{očekivani} broj čvorova ne pređe 30\,000, sa apsolutnim maksimumom od 10 nivoa. Kako se stvarni broj
dece po čvoru izvlači iz Gausove raspodele, a ne fiksira na tu procenu, realizovano stablo
može znatno odstupiti u oba pravca, otuda u konačnom skupu ima stabala i sa svega stotinak i
sa preko 50\,000 čvorova.

\textbf{Algoritmi.}
\begin{itemize}[noitemsep]
    \item \textbf{Minimax} (tačni rešavač): iscrpna rekurzija bez ograničenja budžeta ili dubine.
    Uvek tačan po definiciji; služi isključivo za obeležavanje.
    \item \textbf{Alpha-Beta} (AB): iterativno produbljivanje, pretražuje na dubini 1, pa 2,
    itd., dok budžet čvorova ne bude na izmaku; zadržava se vrednosti poteza sa poslednje
    potpuno završene dubine. AB nigde ne čita tip čvora, nego alternira
    $\max$/$\min$ isključivo po parnosti dubine poziva, pa \textsc{Chance} čvorove tretira kao
    obične \textsc{Max}/\textsc{Min} čvorove.
    \item \textbf{MCTS} (UCT): budžet se meri brojem \textit{dodirnutih} čvorova tokom
    selekcije, ekspanzije i rollout-a do pravog lista, odnosno
    rollout se takođe naplaćuje iz budžeta. Pri selekciji,
    \textsc{Chance} čvorovi se uzorkuju ravnomerno slučajno, dok se
    \textsc{Max}/\textsc{Min} čvorovi biraju UCB1 formulom
    $Q(k) + c\sqrt{\ln N(\text{roditelj})/N(k)}$, $c = \sqrt{2}$, sa $Q$ negiranim na
    \textsc{Min} čvorovima. Finalni potez je dete korena sa najviše poseta.
\end{itemize}

Pošto je MCTS nasumičan, svako stablo se rešava sa 5 nezavisnih, seedovanih
ponavljanja MCTS-a, pa je \texttt{mcts\_correct}
definisan kao većinski glas (bar 3 od 5 tačna), a \texttt{mcts\_correct\_rate} kao kontinualan
udeo tačnih ponavljanja. 

Isti čvorni budžet, izveden kao udeo ukupnog broja čvorova stabla
(log-uniformno iz $[0.02, 0.7]$, sa minimumom od 20), koristi se za AB i za svako MCTS
ponavljanje, čime je poređenje pravedno: oba algoritma rade pod istim ograničenjem resursa.

Generisano je $20\,000$ stabala; zadržano je $16\,978$ (ostatak je odbačen jer je imao manje
od 80 čvorova ili manje od dva deteta korena, dakle degenerisani slučajevi bez stvarnog izbora).

\subsection{Atributi}

Skup ima 91 kolonu: 81 prediktorski atribut i 10 ciljnih/izvedenih promenljivih
(\texttt{which\_algo}, \texttt{which\_better}, \texttt{margin\_category}, \texttt{ab\_correct},
\texttt{mcts\_correct}, \texttt{mcts\_correct\_rate}, \texttt{is\_contested}, \texttt{delta},
\texttt{ab\_regret}, \texttt{mcts\_regret}). Prediktorski atributi se dele u sedam grupa, prema
tome da li se mogu izračunati iz same strukture stabla ili zahtevaju da neki algoritam zaista bude pokrenut:

\begin{itemize}[noitemsep]
    \item \textbf{Parametri generisanja (4):} \texttt{branching\_factor}, \texttt{depth},
    \texttt{chance\_node\_density}, \texttt{heur\_noise\_param};
    \item \textbf{Globalna struktura (24):} brojevi čvorova različitih tipova
    (\texttt{total\_nodes}/\texttt{leaves}, \texttt{internal\_nodes}/\texttt{chance\_nodes}/%
    \texttt{max\_nodes}/\texttt{min\_nodes}),
    izvedeni odnosi 
    
    (\texttt{leaf\_to\_internal\_ratio}, \texttt{chance\_node\_fraction},
    \texttt{leaf\_density}), prosek i standardno odstupanje broja dece po čvoru
    (\texttt{branching\_mean}/\texttt{std}), ostvarena dubina lista (\texttt{realized\_max\_depth},
    \texttt{leaf\_depth\_mean}/\texttt{std}) i deset opisnih statistika vrednosti listova
    (\texttt{leaf\_value\_mean}/\texttt{std}/\texttt{min}/\texttt{max}/\texttt{range}/%
    \texttt{median}, \texttt{leaf\_iqr},
    \item \textbf{Struktura po dubini (25):} \texttt{node\_count\_d\textit{k}},
    \texttt{chance\_node\_count\_d\textit{k}}, 
    
    \texttt{leaf\_value\_mean\_d\textit{k}},
    \texttt{leaf\_value\_std\_d\textit{k}}, \texttt{branching\_mean\_d\textit{k}} za
    $k = 1, \ldots, 5$;
    \item \textbf{Budžet i grananje iz korena (12):} \texttt{node\_budget},
    \texttt{budget\_coverage}, \texttt{ab\_reach\_depth} ($= \log(\mathit{budget}) /
    \log(\mathit{bf})$, koliko nivoa dubine AB realno dostiže), \texttt{depth\_deficit},
    \texttt{depth\_deficit\_rel}, \texttt{branch\_top2\_gap}, \texttt{branch\_mean\_spread},
    \texttt{branch\_value\_skew}, \texttt{top2\_size\_ratio}, \texttt{top2\_depth\_diff},
    \texttt{chance\_frac\_top1/top2}
    \item \textbf{Kvalitet heuristike (3):} \texttt{heur\_incoherence} (i njegovo standardno
    odstupanje \texttt{heur\_incoherence\_std}) meri koliko se heuristička vrednost čvora
    razlikuje od one koju bi dala prava minimax/mean propagacija dece,
    \texttt{heur\_noise\_to\_signal} poredi zadati šum heuristike sa stvarnim rasipanjem
    vrednosti listova. Iako opisuju kvalitet evaluacione funkcije, sva tri atributa se računaju
    direktno iz heurističkih i pravih vrednosti dodeljenih pri generisanju stabla, pa su
    dostupni i pre pokretanja bilo kog algoritma pretrage
    \item \textbf{Atributi ponašanja AB-a i MCTS-a (9):} merenja koja postoje samo zato
    što je algoritam zaista pokrenut: kod AB-a \texttt{ab\_nodes\_visited},
    \texttt{ab\_pruning\_rate}, \texttt{ab\_root\_eval\_std}, \texttt{ab\_instability},
    \texttt{ab\_depth\_reached}, \texttt{ab\_ranking\_correlation}, kod MCTS-a
    \texttt{mcts\_vote\_margin}, \texttt{mcts\_visit\_entropy},
    \texttt{mcts\_best\_value\_estimate}
    \item \textbf{Atributi tačnog rešavača (4):} tačne vrednosti iz iscrpnog minimaxa
    (\texttt{exact\_best\_value}, \texttt{exact\_second\_best\_value},
    \texttt{best\_move\_margin}, \texttt{best\_move\_margin\_normalized}), dostupne tek nakon
    što se Minimax pokrene
\end{itemize}

Skup atributa po dubini pokriva prvih pet nivoa ($k=1,\ldots,5$), umesto svih nivoa do
maksimalne dubine stabla (koja ide i do 10): dublji nivoi bi na manjim stablima gotovo uvek
bili prazni ili nula, pa ne bi nosili stvaran signal, samo bi indirektno kodirali faktor
grananja preko toga koliko kolona uopšte ima nenultu vrednost, što bi kvarilo generalizaciju

\subsection{Ciljne promenljive}

\begin{itemize}[noitemsep]
    \item \texttt{which\_better} (nominalna, 3 klase): čiji je izabrani potez bolji,
    \texttt{AB / MCTS / Tie}; glavni cilj;
    \item \texttt{which\_algo} (nominalna, 4 klase): koji algoritam pronalazi optimalni potez,
    \texttt{AB / MCTS / Both / Neither}; sporedni cilj;
    \item \texttt{ab\_correct}, \texttt{mcts\_correct} (binarne): da li je svaki algoritam
    pojedinačno tačan;
    \item \texttt{mcts\_correct\_rate} (kontinualna): udeo od 5 MCTS ponavljanja koja su
    pogodila optimum;
    \item \texttt{margin\_category} (ordinalna, 3 klase): za koliko pobednik vodi,
    \texttt{low / medium / high};
    \item \texttt{is\_contested} (binarna): identično \texttt{margin\_category = low}
    (prag $0.1$).
\end{itemize}

\texttt{which\_algo} i \texttt{which\_better} odgovaraju na različita pitanja: algoritam koji
promaši baš optimalan potez, ali izabere drugi najbolji, gubi po \texttt{which\_algo} a može
pobediti po \texttt{which\_better}. U pretprocesiranju se iz \texttt{which\_algo} izvodi i
\texttt{tree\_difficulty} (\texttt{Both}->easy, \texttt{AB}/\texttt{MCTS}->medium,
\texttt{Neither}->hard).

\subsection{Raspodela klasa}

\begin{table}[H]
\centering
\caption{Raspodela klasa \texttt{which\_algo} i \texttt{which\_better}}
\begin{tabular}{lrr@{\hspace{2.5em}}lrr}
\toprule
\multicolumn{3}{c}{\texttt{which\_algo}} & \multicolumn{3}{c}{\texttt{which\_better}} \\
\textbf{Klasa} & \textbf{Broj} & \textbf{Udeo} & \textbf{Klasa} & \textbf{Broj} & \textbf{Udeo} \\
\midrule
Both    & 6203 & 36.5\% & AB   & 6330 & 37.3\% \\
Neither & 5445 & 32.1\% & Tie  & 5684 & 33.5\% \\
MCTS    & 2756 & 16.2\% & MCTS & 4964 & 29.2\% \\
AB      & 2574 & 15.2\% &      &      &        \\
\midrule
Ukupno  & 16\,978 & 100\% & Ukupno & 16\,978 & 100\% \\
\bottomrule
\end{tabular}
\end{table}

\texttt{which\_algo} je neravnomerno: klasa Both je više od dva puta zastupljenija od AB
(tabela iznad), jer oba algoritma budu tačna kad god je budžet makar umeren ili je stablo
strukturno lako. Kod \texttt{which\_better} ta neravnoteža je znatno blaža: sve tri klase su
reda veličine 30--37\%, iako AB ima blagu prednost nad MCTS-om. Razlog je verovatno to što
MCTS-ov budžet uključuje i cenu rolauta, pa MCTS pod istim brojem dodirnutih
čvorova stigne da odigra nešto manje selekcije/ekspanzije nego AB. Uprkos
tome, \texttt{which\_better} ostaje uravnoteženiji cilj od \texttt{which\_algo}, što je jedan od razloga zašto je
odabran za glavni cilj klasifikacije.

\begin{figure}[H]
\centering
\includegraphics[width=0.82\textwidth]{../plots/eda/01_class_distribution.png}
\caption{Raspodela klasa \texttt{which\_algo}}
\end{figure}

\begin{figure}[H]
\centering
\includegraphics[width=0.82\textwidth]{../plots/eda/01b_which_better_distribution.png}
\caption{Raspodela klasa \texttt{which\_better}}
\end{figure}

MCTS-ova nestabilnost je izmerena direktno: na 64.6\% stabala se svih 5 ponavljanja
slaže (uvek tačno ili uvek pogrešno), dok na 35.4\% stabala ponavljanja daju različit ishod.

\begin{figure}[H]
\centering
\includegraphics[width=0.75\textwidth]{../plots/eda/03b_mcts_correct_rate.png}
\caption{Udeo od 5 MCTS ponavljanja koja su pogodila optimum, po stablu}
\end{figure}

\subsection{Opisna statistika}

Skup ima 16\,978 slogova i 91 atribut: 62 tipa \texttt{float64}, 26 tipa \texttt{int64} i 3
tekstualna (\texttt{which\_algo}, \texttt{margin\_category}, \texttt{which\_better}). Nema
nedostajućih vrednosti.

\begin{table}[H]
\centering
\caption{Opseg i srednja vrednost parametara generisanja i budžeta}
\begin{tabular}{lrrr}
\toprule
\textbf{Atribut} & \textbf{Min} & \textbf{Srednja} & \textbf{Maks} \\
\midrule
branching\_factor    & 2    & 8.39  & 14  \\
depth                & 2    & 3.97  & 10  \\
chance\_node\_density & 0.0  & 0.38  & 0.9 \\
heur\_noise\_param    & 0.0  & 0.75  & 1.5 \\
budget\_coverage     & 2.0\% & 20.5\% & 69.9\% \\
\bottomrule
\end{tabular}
\end{table}

Gustina CHANCE čvorova ima primetnu masu tačno na nuli, namerna posledica, kako bi imali dodatnih 15\%
potpuno determinističkih stabala koja pogoduju AB-u.

Veličina stabla je izrazito desno zakošena: prosečno 5599 čvorova, ali medijana svega 1472,
jer broj čvorova raste eksponencijalno sa faktorom grananja, pa nekoliko velikih stabala (do
51\,771 čvorova) povlači srednju vrednost znatno iznad tipičnog slučaja.
\texttt{ab\_reach\_depth} (koliko nivoa dubine budžet realno dostiže) prati istu logiku u
suprotnom smeru: prosečno dostiže tek 2.88 nivoa.

\begin{figure}[H]
\centering
\includegraphics[width=0.9\textwidth]{../plots/eda/04_parameter_distribution.png}
\caption{Raspodela parametara generisanja i budžeta}
\end{figure}

%-------------------------------------------------------------------
\section{Eksplorativna analiza podataka}
%-------------------------------------------------------------------

\subsection{Atributi po klasama}

\begin{figure}[H]
\centering
\includegraphics[width=0.95\textwidth]{../plots/eda/05_feature_boxplots.png}
\caption{Raspodela šest ključnih atributa po klasi \texttt{which\_algo}}
\end{figure}

Faktor grananja je najizrazitiji strukturni separator: stabla iz klase Neither imaju u
proseku najveći faktor grananja (9.28), Both najmanji (7.42), dok AB i MCTS leže između
njih dvoje i liče jedna na drugu (slika iznad). Gustina CHANCE čvorova pre svega izdvaja
MCTS, koja je u proseku gušća nego kod preostale tri klase.

\subsection{Korelaciona analiza}

\begin{figure}[H]
\centering
\includegraphics[width=0.95\textwidth]{../plots/eda/07_ab_mcts_correlation.png}
\caption{Pearsonova korelacija numeričkih atributa sa \texttt{ab\_correct} i \texttt{mcts\_correct}}
\end{figure}

\begin{figure}[H]
\centering
\includegraphics[width=0.9\textwidth]{../plots/eda/08_correlation_matrix.png}
\caption{Korelaciona matrica odabranih atributa}
\end{figure}

\begin{table}[H]
\centering
\caption{Najjače Pearsonove korelacije sa \texttt{ab\_correct} i \texttt{mcts\_correct}}
\begin{tabular}{lrr}
\toprule
\textbf{Atribut} & \textbf{ab\_correct} & \textbf{mcts\_correct} \\
\midrule
ab\_ranking\_correlation & 0.398  & --     \\
mcts\_vote\_margin       & 0.286  & 0.441  \\
branch\_top2\_gap        & 0.257  & --     \\
ab\_reach\_depth         & 0.179  & 0.357  \\
depth                    & 0.175  & 0.316  \\
best\_move\_margin       & 0.240  & --     \\
branching\_factor        & $-0.175$ & $-0.198$ \\
mcts\_visit\_entropy     & $-0.265$ & $-0.382$ \\
\bottomrule
\end{tabular}
\end{table}

Za oba cilja, najjača korelacija dolazi od atributa ponašanja algoritama
(\texttt{ab\_ranking\_correlation}, \texttt{mcts\_vote\_margin}), što je očekivano jer oba
direktno mere sigurnost algoritma u trenutku odluke, ne samo opštu težinu stabla. Od
strukturnih atributa, dostupnih pre pokretanja, budžet i margina korena
(\texttt{branch\_top2\_gap}, \texttt{best\_move\_margin}) pokazuju najjaču pozitivnu vezu, a
veći faktor grananja ide uz nižu tačnost oba algoritma.

Zanimljivo je da \texttt{ab\_reach\_depth} i \texttt{depth} jako koreliraju i sa
\texttt{mcts\_correct}, iako direktno mere AB-ov budžet, ne MCTS-ov. Razlog je da veći
\texttt{depth} u ovom skupu ide uz manji faktor grananja (stabla su ograničena na
$\approx 30\,000$ čvorova), pa oba algoritma profitiraju od istog, lakšeg tipa stabla.
\texttt{mcts\_visit\_entropy} je negativno korelisan sa oba cilja: visoka entropija poseta
(MCTS ravnomerno raspoređuje posete po svoj deci umesto da se skoncentriše na jedno) je
opštiji signal da je stablo teško, ne samo da je MCTS neodlučan.

\subsection{Vizualizacija u 2D prostoru}

\textbf{PCA.} Primenjena na 68 strukturnih atributa (izuzimajući ciljne kolone i kolone koje
opisuju ponašanje algoritama). Prve dve komponente objašnjavaju zajedno 34.7\% varijanse
(PC1: 21.9\%, PC2: 12.8\%),
a potrebno je 13 komponenti za 80\% varijanse. 

\begin{figure}[H]
\centering
\begin{subfigure}[b]{0.54\textwidth}
    \includegraphics[width=\textwidth]{../plots/eda/09_pca_scatter.png}
    \caption{PCA rasejani dijagram, obojen po \texttt{which\_algo}}
\end{subfigure}
\hfill
\begin{subfigure}[b]{0.43\textwidth}
    \includegraphics[width=\textwidth]{../plots/eda/10_pca_variance.png}
    \caption{Kumulativna objašnjena varijansa}
\end{subfigure}
\caption{PCA vizualizacija}
\end{figure}

Klase nisu linearno razdvojive u 2D projekciji: gotovo 65\% informacije se gubi pri
svođenju na dve komponente, pa je za pouzdanu klasifikaciju potreban nelinearan model
(odeljak 5).

\textbf{t-SNE.} Nelinearna vizualizacija na prvih 26 PCA
komponenti, uzorak od 4000 stabala.

\begin{figure}[H]
\centering
\includegraphics[width=0.7\textwidth]{../plots/eda/11_tsne.png}
\caption{t-SNE vizualizacija, obojena po \texttt{which\_algo}}
\end{figure}

t-SNE otkriva lokalne nakupine pretežno jedne klase, ali sa značajnim preklapanjem. Problem je,
dakle, naučiv iz strukturnih atributa, ali nije trivijalan.

%-------------------------------------------------------------------
\section{Pretprocesiranje podataka}
%-------------------------------------------------------------------

\subsection{Izvedeni atributi}

\begin{itemize}[noitemsep]
    \item \texttt{ab\_coverage} $= \mathtt{ab\_nodes\_visited} / \mathtt{total\_nodes}$:
    koliko je AB stvarno obišao, za razliku od \texttt{budget\_coverage} koji je
    \textit{nameravani} budžet. Srednja vrednost 0.201, korelacija sa \texttt{budget\_coverage}
    0.993 (skoro identične, jer AB retko iscrpi budžet pre nego što dostigne kraj iteracije);
    \item \texttt{mcts\_unanimous}: da li se svih 5 MCTS ponavljanja slaže (64.6\%, videti
    odeljak 2.4);
    \item \texttt{depth\_chance\_inter} $= \mathtt{depth} \times \mathtt{chance\_node\_density}$:
    jeftina interakcija dve ose generatora;
    \item \texttt{is\_tie}: \texttt{which\_better = Tie} (5684 od 16\,978, 33.5\%);
    \item \texttt{tree\_difficulty}: easy/medium/hard izvedeno iz \texttt{which\_algo}
    (6203 / 5330 / 5445).
\end{itemize}

Ove poslednje dve, \texttt{ab\_coverage} i \texttt{mcts\_unanimous}, spadaju u atribute
ponašanja algoritama (odeljak 5.1 ih izbacuje iz skupa B), jer zahtevaju da AB, odnosno MCTS,
zaista budu
pokrenuti.

\subsection{Čišćenje podataka}

Skup nema nedostajućih vrednosti. Autlajeri su provereni IQR i $3\sigma$ metodom na šest
ključnih atributa:

\begin{table}[H]
\centering
\caption{Analiza autlajera}
\begin{tabular}{lrrrr}
\toprule
\textbf{Atribut} & \textbf{IQR autlajeri} & \textbf{\%} & \textbf{$3\sigma$ autlajeri} & \textbf{\%} \\
\midrule
total\_nodes       & 1748 & 10.3\% & 364 & 2.1\% \\
best\_move\_margin & 855  & 5.0\%  & 328 & 1.9\% \\
mcts\_vote\_margin & 1533 & 9.0\%  & 346 & 2.0\% \\
heur\_incoherence  & 403  & 2.4\%  & 197 & 1.2\% \\
leaf\_value\_std   & 66   & 0.4\%  & 68  & 0.4\% \\
ab\_pruning\_rate  & 2    & 0.0\%  & 11  & 0.1\% \\
\bottomrule
\end{tabular}
\end{table}

Autlajeri se zadržavaju: nisu greške merenja, već realne ali ređe konfiguracije
stabla (npr. veoma veliko stablo je validan, ne pogrešan, parametar). Njihovo brisanje bi
uklonilo informaciju relevantnu za klasifikaciju.

\subsection{Kodiranje kategoričkih atributa}

\texttt{margin\_category} i \texttt{tree\_difficulty} su ordinalne i kodiraju se
\texttt{OrdinalEncoder}-om sa eksplicitnim redosledom (\texttt{low<medium<high},
\texttt{easy<medium<hard}). \texttt{which\_algo} i \texttt{which\_better} su nominalne i
kodiraju se \texttt{LabelEncoder}-om: \texttt{which\_algo}
$-> \{\text{AB}{=}0, \text{Both}{=}1, \text{MCTS}{=}2, \text{Neither}{=}3\}$,
\texttt{which\_better} $-> \{\text{AB}{=}0, \text{MCTS}{=}1, \text{Tie}{=}2\}$.

\subsection{Standardizacija}

84 numerička prediktorska atributa (sve osim ciljnih/izvedenih kolona) se standardizuju
\texttt{StandardScaler}-om, $z = (x-\mu)/\sigma$, obavezno pre logističke regresije i
neuronske mreže. Naučeni scaler se čuva u \texttt{models/scaler.pkl} radi identične
transformacije na budućim podacima bez ponovnog fitovanja.

\subsection{Izbor podskupa atributa}

Cilj izbora je \texttt{which\_better}. Tri grupe metoda glasaju nezavisno:

\textbf{Filter.} Korelaciona eliminacija ($|r| \geq 0.99$): $84 -> 74$ atributa
(\texttt{total\_leaves} $\leftrightarrow$ \texttt{total\_nodes} $r{=}0.995$;
\texttt{leaf\_to\_internal\_ratio} $\leftrightarrow$ \texttt{branching\_factor} $r{=}0.994$;
\texttt{branching\_mean} $\leftrightarrow$ \texttt{branching\_factor} $r{=}0.994$;
\texttt{realized\_max\_depth} $\leftrightarrow$ \texttt{depth} $r{=}1.000$;
\texttt{leaf\_depth\_mean} $\leftrightarrow$ \texttt{depth} $r{=}1.000$;
\texttt{leaf\_fraction\_positive} $\leftrightarrow$ \texttt{leaf\_value\_mean} $r{=}0.991$;
\texttt{branching\_mean\_d1} $\leftrightarrow$ \texttt{leaf\_to\_internal\_ratio} $r{=}0.993$;
\texttt{ab\_nodes\_visited} $\leftrightarrow$ \texttt{node\_budget} $r{=}1.000$;
\texttt{best\_move\_margin\_normalized} $\leftrightarrow$ \texttt{best\_move\_margin}
$r{=}0.998$; \texttt{ab\_coverage} $\leftrightarrow$ \texttt{budget\_coverage} $r{=}0.993$).
Kvazi-konstantni prag (jedna vrednost $>99\%$ redova) ne eliminiše ništa ($74 -> 74$). ANOVA
F-test (\texttt{SelectKBest}) bira top 37 od 74.

\textbf{Wrapper.} Rekurzivna eliminacija (RFE, logistička regresija) do $74 -> 37$ atributa.

\textbf{Embedded.} Slučajne šume (100 stabala), top 37 po Ginijevoj važnosti.

\begin{figure}[H]
\centering
\includegraphics[width=0.7\textwidth]{../plots/preprocessing/rf_importance.png}
\caption{Top 20 atributa po važnosti u Slučajnim šumama (izbor atributa)}
\end{figure}

\textbf{Ansambl glasanje.} Atribut ulazi u finalni skup ako ga izabere bar 2 od 3 grupe:

\begin{table}[H]
\centering
\caption{Finalni skup atributa (33 atributa, $\geq 2/3$ glasova)}
\small
\begin{tabular}{lcccc}
\toprule
\textbf{Atribut} & \textbf{ANOVA} & \textbf{RFE} & \textbf{RF} & \textbf{Glasovi} \\
\midrule
ab\_depth\_reached       & X & X & X & 3 \\
ab\_instability          & X & X & X & 3 \\
ab\_pruning\_rate        & X & X & X & 3 \\
ab\_reach\_depth         & X & X & X & 3 \\
best\_move\_margin       & X & X & X & 3 \\
branch\_top2\_gap        & X & X & X & 3 \\
budget\_coverage         & X & X & X & 3 \\
depth\_deficit\_rel      & X & X & X & 3 \\
leaf\_depth\_std         & X & X & X & 3 \\
mcts\_best\_value\_estimate & X & X & X & 3 \\
mcts\_unanimous          & X & X & X & 3 \\
mcts\_visit\_entropy     & X & X & X & 3 \\
mcts\_vote\_margin       & X & X & X & 3 \\
node\_budget             & X & X & X & 3 \\
total\_internal\_nodes   & X & X & X & 3 \\
total\_nodes             & X & X & X & 3 \\
ab\_root\_eval\_std      &   & X & X & 2 \\
branch\_mean\_spread     &   & X & X & 2 \\
branch\_value\_skew      & X &   & X & 2 \\
branching\_mean\_d2      & X & X &   & 2 \\
branching\_mean\_d4      & X & X &   & 2 \\
depth                    & X & X &   & 2 \\
depth\_chance\_inter     & X & X &   & 2 \\
depth\_deficit           &   & X & X & 2 \\
exact\_second\_best\_value &  & X & X & 2 \\
heur\_incoherence        & X &   & X & 2 \\
leaf\_density            &   & X & X & 2 \\
leaf\_value\_range       & X &   & X & 2 \\
leaf\_value\_std\_d3     & X & X &   & 2 \\
leaf\_value\_std\_d5     & X & X &   & 2 \\
node\_count\_d4          & X & X &   & 2 \\
total\_min\_nodes        & X &   & X & 2 \\
\bottomrule
\end{tabular}
\end{table}

Finalni pretprocesirani skup (\texttt{data/trees\_preprocessed.csv}): 16\,978 slogova, 33
atributa (\texttt{feat\_A}) plus 16 ciljnih/izvedenih kolona (49 kolona ukupno).

%-------------------------------------------------------------------
\section{Klasifikacija}
%-------------------------------------------------------------------

\subsection{Pregled eksperimenta}

Glavni cilj je \texttt{which\_better} (3 klase); sporedni cilj, obrađen paralelno, je
\texttt{which\_algo} (odeljak 5.6), uz dodatne modele za \texttt{margin\_category} i
\texttt{tree\_difficulty} (odeljak 5.7) i binarne \texttt{ab\_correct}/\texttt{mcts\_correct}
(odeljak 5.5).

Eksperiment radimo na tri skupa atributa:

\begin{itemize}[noitemsep]
    \item \textbf{Skup A (33 atributa):} finalni pretprocesirani skup
    \item \textbf{Skup B (28 atributa):} samo strukturni atributi, bez ijednog merenja koje
    zahteva da se algoritmi zaista pokrenu (\texttt{ab\_nodes\_visited},
    \texttt{mcts\_vote\_margin},
    itd.); \texttt{heur\_incoherence} i \texttt{ab\_reach\_depth} ostaju, jer se računaju iz
    same strukture. Skup se dobija tako što se iz skupa A ukloni 11 atributa ponašanja
    algoritama (22 preostala atributa), a zatim mu se vrati 6 dodatnih strukturnih atributa
    koji nisu prošli izbor u odeljku 4.6 (\texttt{leaf\_to\_internal\_ratio},
    \texttt{node\_count\_d3}, \texttt{leaf\_value\_std\_d2}, \texttt{chance\_node\_count\_d2},
    \texttt{total\_max\_nodes}, \texttt{heur\_noise\_param});
    \item \textbf{Skup C (81 atribut):} svi prediktorski atributi originalnog skupa, bez
    izbora.
\end{itemize}

Skup B odgovara na centralno pitanje: \textit{može li se, pre pokretanja ijednog algoritma,
iz same strukture stabla predvideti čiji će potez biti bolji?} Razlika B$->$A meri vrednost
merenja ponašanja algoritama; razlika A$->$C meri doprinos izbora atributa.

\textbf{Protokol.} Stratifikovana unakrsna validacija, $k{=}10$. Primarna metrika je
macro-F1: harmonijska sredina preciznosti i odziva, izračunata posebno za svaku klasu i
zatim neponderisano usrednjena, tako da model ne može da izgleda dobro samo time što
pogađa najčešću klasu. Ovo je bitno jer \texttt{which\_algo} ostaje neravnomeran (odeljak
2.4). Dodatno se prati i ROC AUC (jedan-naspram-ostalih, OvR): meri koliko dobro model
rangira instance jedne klase iznad ostalih, nezavisno od izabranog praga odlučivanja, gde
$0.5$ odgovara slučajnom nagađanju, a $1.0$ savršenom razdvajanju.

\textbf{Korišćeni algoritmi.}
\begin{itemize}[noitemsep]
    \item \textbf{Stablo odlučivanja (DT):} rekurzivno deli prostor atributa pragovima na
    pojedinačnim atributima. Lako se vizualizuje i čita direktno (odeljak 5.3), pa služi kao
    interpretabilna polazna tačka pre složenijih modela.
    \item \textbf{Slučajne šume (RF, 150 stabala):} ansambl stabala odlučivanja treniranih
    nad slučajnim poduzorcima podataka i atributa, čije se predikcije usrednjavaju
    glasanjem. Smanjuje varijansu pojedinačnog stabla i obično daje solidan rezultat bez
    detaljnog podešavanja hiperparametara.
    \item \textbf{Logistička regresija (LR):} linearni probabilistički klasifikator.
    Uključena je kao provera da li linearna granica između klasa uopšte može da objasni
    razliku.
    \item \textbf{XGBoost (XGB, 150 stabala):} gradijentno bustovan ansambl stabala, gde
    svako naredno stablo ispravlja greške prethodnih. Ona u praksi daje još bolje rezultate od prethodnih i predstavlja jači ansambl metod u poređenju.
    \item \textbf{Neuronska mreža (MLP, jedan skriveni sloj od 64 neurona):} uči nelinearne
    kombinacije atributa gradijentnim spustom. Testira da li dodatna nelinearnost donosi
    prednost nad modelima baziranim na stablima.
    \item \textbf{Metoda potpornih vektora (SVM, RBF kernel, $C{=}1.0$):} traži granicu
    odlučivanja koja maksimizuje marginu između klasa u (kernelom transformisanom) prostoru
    atributa.
    \item \textbf{Slaganje modela (Stacking: DT+RF+XGB $->$ logistička regresija} kombinuje predikcije prethodna tri modela kao ulaz za finalnu logističku
    regresiju, koja uči kako da ih najbolje izmeša. Ideja je da modeli greše na različitim
    primerima, pa njihova kombinacija može biti robusnija od bilo kog modela pojedinačno.
\end{itemize}

\subsection{Rezultati: which\_better}

\begin{table}[H]
\centering
\caption{Macro-F1 po modelu i skupu atributa (10-fold, stratifikovano)}
\begin{tabular}{lccc}
\toprule
\textbf{Model} & \textbf{A (33 atr.)} & \textbf{B (28 atr.)} & \textbf{C (81 atr.)} \\
\midrule
Stablo odlučivanja (DT)   & 0.622 & 0.470 & 0.564 \\
Slučajne šume (RF)        & 0.679 & 0.554 & 0.649 \\
Logistička regresija (LR) & 0.679 & 0.557 & 0.652 \\
XGBoost (XGB)              & 0.689 & 0.556 & 0.664 \\
Neuronska mreža (MLP)      & 0.689 & 0.561 & 0.660 \\
SVM                        & 0.682 & 0.554 & 0.656 \\
Slaganje modela (Stacking) & \textbf{0.692} & --- & --- \\
\bottomrule
\end{tabular}
\end{table}

Slaganje modela na skupu A postiže najviši macro-F1 (0.692), blago iznad XGBoost-a i MLP-a
(oba 0.689). Razlike među modelima na skupu A su male, što govori da je informacija u atributima
već zasićena: izbor modela tu dodaje malo iznad izbora atributa. Pad sa A na B
(0.679 $->$ 0.554 za RF) je znatan: merenja ponašanja algoritama nose realnu dodatnu
vrednost, ali skup B sam po sebi, bez pokretanja algoritama, već daje macro-F1 blizu 0.55,
znatno iznad slučajnog pogađanja za tri klase ($\approx 0.33$).

\begin{table}[H]
\centering
\caption{Detaljne mere na skupu A (\texttt{which\_better}, 10-fold)}
\begin{tabular}{lrrrrr}
\toprule
\textbf{Model} & \textbf{Tačnost} & \textbf{macro-F1} & \textbf{Preciznost} & \textbf{Odziv} & \textbf{ROC AUC (OvR)} \\
\midrule
DT       & 0.629 & 0.622 & 0.622 & 0.622 & 0.719 \\
RF       & 0.701 & 0.679 & 0.686 & 0.686 & 0.873 \\
LR       & 0.697 & 0.679 & 0.683 & 0.684 & 0.870 \\
XGB      & 0.701 & 0.689 & 0.689 & 0.691 & 0.876 \\
MLP      & 0.705 & 0.690 & 0.691 & 0.694 & 0.876 \\
SVM      & 0.706 & 0.682 & 0.692 & 0.691 & 0.872 \\
Stacking & \textbf{0.707} & \textbf{0.692} & 0.693 & \textbf{0.695} & \textbf{0.876} \\
\bottomrule
\end{tabular}
\end{table}

\begin{figure}[H]
\centering
\includegraphics[width=0.95\textwidth]{../plots/classification/01_confusion_matrices.png}
\caption{Matrice konfuzije, svi modeli, \texttt{which\_better}, skup A}
\end{figure}

Kada se klasa Tie u potpunosti isključi (stabla na kojima AB i MCTS biraju potez iste
vrednosti), problem postaje binarna klasifikacija AB naspram MCTS na preostalih 11\,294
stabala (66.5\% skupa). Macro-F1 raste kod svih modela, od DT-a ($0.641$) do Stacking-a
($0.713$), koji ostaje najbolji i ovde. Isključivanje Tie-a, dakle, čini problem lakšim za sve
modele podjednako, bez promene njihovog relativnog poretka.

\begin{figure}[H]
\centering
\includegraphics[width=0.95\textwidth]{../plots/classification/12_confusion_matrices_no_tie.png}
\caption{Matrice konfuzije, svi modeli, \texttt{which\_better} bez klase Tie, skup A}
\end{figure}

\subsection{Interpretabilni model: stablo odlučivanja}

Stablo odlučivanja (\texttt{max\_depth=5}) trenirano na skupu B postiže tačnost 0.589 na
trening skupu i deli koren po \texttt{node\_budget} $\leq 130.5$. Prvi i najvažniji signal je,
dakle, da li AB uopšte ima dovoljno apsolutnog budžeta na raspolaganju. Drugi po važnosti signal, dalje u
stablu, je \texttt{branch\_top2\_gap}, strukturna margina između dva najbolja poteza iz
korena.

\begin{figure}[H]
\centering
\includegraphics[width=0.95\textwidth]{../plots/classification/02_decision_tree.png}
\caption{Stablo odlučivanja (\texttt{max\_depth=5}, skup B, \texttt{which\_better})}
\end{figure}

\subsection{Važnost atributa}

\begin{table}[H]
\centering
\caption{Top 10 atributa po važnosti (skup B, \texttt{which\_better})}
\begin{tabular}{lr@{\hspace{3em}}lr}
\toprule
\multicolumn{2}{c}{\textbf{Slučajne šume (Gini)}} & \multicolumn{2}{c}{\textbf{XGBoost (Gain)}} \\
\midrule
branch\_top2\_gap       & 0.093 & node\_budget            & 0.186 \\
node\_budget            & 0.068 & branch\_top2\_gap       & 0.074 \\
ab\_reach\_depth        & 0.052 & chance\_node\_count\_d2 & 0.039 \\
branch\_value\_skew     & 0.050 & ab\_reach\_depth        & 0.037 \\
branch\_mean\_spread    & 0.047 & budget\_coverage        & 0.034 \\
heur\_noise\_param      & 0.045 & branch\_mean\_spread    & 0.032 \\
budget\_coverage        & 0.045 & heur\_noise\_param      & 0.031 \\
depth\_deficit\_rel     & 0.043 & total\_nodes            & 0.031 \\
total\_nodes            & 0.041 & total\_internal\_nodes  & 0.031 \\
heur\_incoherence       & 0.040 & total\_max\_nodes       & 0.030 \\
\bottomrule
\end{tabular}
\end{table}

Oba modela slažu se da su budžet (\texttt{node\_budget}, \texttt{budget\_coverage}) i strukturna
margina korena (\texttt{branch\_top2\_gap}) najvažniji atributi, u skladu sa
stablom odlučivanja iz odeljka 5.3. Slažu se i da je \texttt{ab\_reach\_depth} relevantan, ali XGBoost mu dodeljuje znatno manju težinu u odnosu na sam \texttt{node\_budget}
(0.186 naspram 0.037), dok Slučajne šume najviše rangiraju \texttt{branch\_top2\_gap} (0.093).

\subsection{Binarna klasifikacija: ab\_correct i mcts\_correct}

\begin{figure}[H]
\centering
\includegraphics[width=0.88\textwidth]{../plots/classification/05_roc_curves.png}
\caption{ROC krive za \texttt{ab\_correct} (levo) i \texttt{mcts\_correct} (desno), RF, skup A naspram skupa B, 5-fold}
\end{figure}

\begin{table}[H]
\centering
\caption{ROC AUC (RF, 5-fold CV, skup A naspram skupa B)}
\begin{tabular}{lrr}
\toprule
\textbf{Cilj} & \textbf{Skup A} & \textbf{Skup B} \\
\midrule
\texttt{ab\_correct}   & 0.739 & 0.696 \\
\texttt{mcts\_correct} & 0.833 & 0.769 \\
\bottomrule
\end{tabular}
\end{table}

Razlika između skupa A i skupa B je malo veća za \texttt{mcts\_correct} nego za
\texttt{mcts\_correct} (tabela iznad). Atributi ponašanja algoritama, dakle, nose relativno
više dodatne informacije o AB-ovoj tačnosti nego o MCTS-ovoj. Ipak, u oba slučaja skup B, koji
ne sadrži nijedno merenje ponašanja algoritama, sam po sebi nosi znatan signal, daleko iznad
AUC $= 0.5$ slučajnog modela. Atributi u skupu A (posebno \texttt{mcts\_vote\_margin}) dodaju realnu vrednost kod oba targeta,
što je očekivano jer ti atributi direktno mere sigurnost algoritma u trenutku odluke.

\subsubsection{SHAP analiza greške AB algoritma}

XGBoost na skupu B, cilj \texttt{ab\_correct}: SHAP (SHapley Additive exPlanations)
vrednosti su zasnovane na Šeplijevim vrednostima iz teorije kooperativnih igara. Pozitivna vrednost gura predikciju ka \texttt{ab\_correct}
$=1$ (AB tačan), negativna ka $0$ (AB pogrešan).

\begin{figure}[H]
\centering
\includegraphics[width=0.95\textwidth]{../plots/classification/06_shap_ab_correct.png}
\caption{SHAP analiza za \texttt{ab\_correct} na skupu B (XGBoost)}
\end{figure}

\begin{table}[H]
\centering
\caption{SHAP doprinosi za \texttt{ab\_correct} (skup B, XGBoost)}
\begin{tabular}{lr@{\hspace{3em}}lr}
\toprule
\multicolumn{2}{c}{\textbf{Pozitivan doprinos (AB tačan)}} & \multicolumn{2}{c}{\textbf{Negativan doprinos (AB pogrešan)}} \\
\midrule
node\_count\_d3         &  0.091 & branching\_mean\_d2   & $-0.077$ \\
ab\_reach\_depth        &  0.039 & total\_min\_nodes     & $-0.038$ \\
branch\_top2\_gap       &  0.039 & total\_internal\_nodes & $-0.034$ \\
total\_nodes            &  0.035 & leaf\_density         & $-0.018$ \\
total\_max\_nodes       &  0.032 & depth\_deficit        & $-0.015$ \\
node\_count\_d4         &  0.020 & node\_budget          & $-0.015$ \\
heur\_incoherence       &  0.018 & leaf\_value\_std\_d3  & $-0.008$ \\
\bottomrule
\end{tabular}
\end{table}

Najveći pojedinačni doprinos ima \texttt{node\_count\_d3} (pozitivan, broj čvorova na trećem
nivou dubine) naspram \texttt{branching\_mean\_d2} (negativan, prosečan broj dece po čvoru na
drugom nivou). Oba su strukturni atributi vezani za oblik stabla blizu korena, ne za budžet
direktno. \texttt{total\_nodes} doprinosi pozitivno (veći ukupan broj čvorova), a
\texttt{node\_budget} negativno (manji budžet). Zaključak je da oblik stabla u
ranim nivoima i budžet zajedno nose najviše signala.
\texttt{branch\_top2\_gap} (strukturna margina) i \texttt{heur\_incoherence} (kvalitet
heuristike) potvrđuju istu poruku kao stablo odlučivanja i tabela važnosti.

\subsection{which\_algo naspram which\_better}

\begin{table}[H]
\centering
\caption{\texttt{which\_algo} (bez klase Neither), RF/XGB, skup B, 10-fold}
\begin{tabular}{lrr}
\toprule
\textbf{Model} & \textbf{macro-F1} & \textbf{std} \\
\midrule
RF  & 0.521 & 0.018 \\
XGB & 0.528 & 0.014 \\
\bottomrule
\end{tabular}
\end{table}

Klasa Neither (5445 stabala, 32.1\%) je isključena, jer ta stabla nemaju zajednički
strukturni potpis, nego samo dodaju šum bez korisnog signala za preostale tri klase. Izveštaj
po klasi (RF, skup B):

\begin{table}[H]
\centering
\begin{tabular}{lrrrr}
\toprule
\textbf{Klasa} & \textbf{Preciznost} & \textbf{Odziv} & \textbf{F1} & \textbf{Broj} \\
\midrule
AB   & 0.61 & 0.52 & 0.56 & 2574 \\
Both & 0.65 & 0.85 & 0.74 & 6203 \\
MCTS & 0.43 & 0.19 & 0.27 & 2756 \\
\bottomrule
\end{tabular}
\end{table}

Poređeno sa glavnim ciljem: \texttt{which\_better} na istom skupu B postiže RF macro-F1
$= 0.554$, što je iznad RF rezultata za \texttt{which\_algo} bez klase Neither (0.521). Odbacivanje
klase Neither pomaže, ali ne dovoljno da \texttt{which\_algo} stigne \texttt{which\_better},
što je dodatna potvrda da je pitanje čiji je potez bolji strukturno čistiji problem od pitanja
ko je pogodio tačno (odeljak 2.4).

\begin{figure}[H]
\centering
\includegraphics[width=0.7\textwidth]{../plots/classification/09_which_algo_confusion.png}
\caption{Matrica konfuzije: \texttt{which\_algo} bez Neither (RF, skup B, 10-fold)}
\end{figure}

Kada se, pored Neither, odbaci i klasa Both (ostaju samo stabla na kojima je tačno jedan od
dva algoritma pogodio optimum), \texttt{which\_algo} postaje binarna klasifikacija AB naspram
MCTS. Macro-F1 tada znatno raste: RF $= 0.712 \pm 0.019$, XGB $= 0.689 \pm 0.015$, daleko iznad
rezultata sa sve tri klase. Izveštaj po klasi (RF, skup B): AB 0.72/0.67/0.69 (2574 slogova),
MCTS 0.71/0.75/0.73 (2756 slogova), ukupna tačnost 0.71. Najveći deo teškoće u
\texttt{which\_algo} dolazi, dakle, od razdvajanja Both/Neither od preostalih klasa, a ne od
razlikovanja same AB-ove i MCTS-ove greške.

\begin{figure}[H]
\centering
\includegraphics[width=0.7\textwidth]{../plots/classification/13_which_algo_ab_vs_mcts.png}
\caption{Matrica konfuzije: \texttt{which\_algo}, samo AB naspram MCTS (RF, skup B, 10-fold)}
\end{figure}

\subsection{Ordinalni ciljevi: margin\_category i tree\_difficulty}

Oba se klasifikuju Slučajnim šumama na skupu B, 10-fold CV.

\begin{table}[H]
\centering
\caption{\texttt{margin\_category} (RF, skup B, 10-fold), macro-F1 $= 0.53$}
\begin{tabular}{lrrrr}
\toprule
\textbf{Klasa} & \textbf{Preciznost} & \textbf{Odziv} & \textbf{F1} & \textbf{Broj} \\
\midrule
low    & 0.60 & 0.46 & 0.52 & 5842 \\
medium & 0.53 & 0.70 & 0.61 & 7639 \\
high   & 0.58 & 0.40 & 0.47 & 3497 \\
\bottomrule
\end{tabular}
\end{table}

\begin{table}[H]
\centering
\caption{\texttt{tree\_difficulty} (RF, skup B, 10-fold), macro-F1 $= 0.49$}
\begin{tabular}{lrrrr}
\toprule
\textbf{Klasa} & \textbf{Preciznost} & \textbf{Odziv} & \textbf{F1} & \textbf{Broj} \\
\midrule
easy   & 0.58 & 0.70 & 0.63 & 6203 \\
medium & 0.41 & 0.32 & 0.36 & 5330 \\
hard   & 0.50 & 0.49 & 0.49 & 5445 \\
\bottomrule
\end{tabular}
\end{table}

Oba modela imaju sličan, umeren macro-F1 ($\approx 0.5$): margina i težina stabla su
delimično predvidive iz strukture. Medium kategorija je
najteže razdvojiva u oba slučaja, jer po definiciji leži na granici prema oba suseda.

%-------------------------------------------------------------------
\section{Vizualizacije i generalizacija}
%-------------------------------------------------------------------

\subsection{Dijagram faza}

Dominantna klasa \texttt{which\_better} u prostoru (\texttt{ab\_reach\_depth} $\times$
gustina CHANCE čvorova, oba binovana u 6 intervala):

AB dominira ceo nisko-budžetni opseg bez obzira na gustinu CHANCE čvorova, jer MCTS sa
malim brojem poseta po grani ne stigne da izgradi pouzdanu statistiku. Tie
dominira ceo visoko budžetni opseg, jer oba algoritma tada uglavnom nađu isti ili blizak
potez. MCTS izranja kao srednja traka, najizraženija baš pri višoj gustini CHANCE čvorova,
gde usrednjavanje po slučajnosti otežava AB-ovu plitku heurističku procenu više nego
MCTS-ovo uzorkovanje.

\subsection{Preciznost rangiranja AB algoritma}

\begin{figure}[H]
\centering
\includegraphics[width=0.7\textwidth]{../plots/visualizations/03_ab_ranking.png}
\caption{Spirmanova korelacija AB-ovog rangiranja poteza po klasi \texttt{which\_algo}}
\end{figure}

Binarna metrika \texttt{ab\_correct} ne razlikuje algoritam koji bira drugi po kvalitetu
potez od onog koji bira najgori. Spirmanova korelacija rangova (\texttt{ab\_ranking\_correlation},
$+1$ = savršeno uređivanje) pokazuje bogatiju sliku po klasi \texttt{which\_algo} kod klasa gde je AB tačan (Both, AB) korelacija je visoka, preko 0.7, dok kod
klasa gde formalno greši (MCTS, Neither) ostaje pozitivna, samo osetno niža, oko 0.45. AB
retko pravi katastrofalnu grešku u rangiranju čak i kada formalno pogreši: identifikuje skup
dobrih poteza, ali ih pogrešno rangira na margini.

\subsection{Generalizacija: predikcija na neviđenom faktoru grananja}

Referentna tačnost, 5-fold CV na celom skupu, skup B: \texttt{which\_algo} macro-F1
$= 0.402 \pm 0.005$, \texttt{which\_better} macro-F1 $= 0.559 \pm 0.008$.

Stroži test: trening bez svih stabala određenog \texttt{branching\_factor}, test isključivo
na njima. \texttt{bf}$=7$ leži usred opsega $[2,14]$ (interpolacija), \texttt{bf}$=14$ je
ivica opsega (ekstrapolacija).

\begin{table}[H]
\centering
\caption{Generalizacija na neviđen \texttt{branching\_factor} (RF, skup B)}
\begin{tabular}{lrrrrrrrr}
\toprule
& \multicolumn{2}{c}{\texttt{ab\_correct}} & \multicolumn{2}{c}{\texttt{mcts\_correct}} & \multicolumn{2}{c}{\texttt{which\_algo}} & \multicolumn{2}{c}{\texttt{which\_better}} \\
\textbf{Scenario} & Tač. & F1 & Tač. & F1 & Tač. & F1 & Tač. & F1 \\
\midrule
bf=7 (interpolacija)   & 0.658 & 0.674 & 0.687 & 0.735 & 0.487 & 0.360 & 0.569 & 0.538 \\
bf=14 (ekstrapolacija) & 0.667 & 0.574 & 0.722 & 0.662 & 0.548 & 0.481 & 0.622 & 0.610 \\
bf=7 i bf=14 (oba)     & 0.653 & 0.637 & 0.704 & 0.717 & 0.506 & 0.400 & 0.590 & 0.567 \\
\bottomrule
\end{tabular}
\end{table}

\begin{figure}[H]
\centering
\includegraphics[width=0.95\textwidth]{../plots/visualizations/04b_generalization_cm_better.png}
\caption{\texttt{which\_better}: predikcija iz strukturnih atributa na neviđenim vrednostima \texttt{branching\_factor}}
\end{figure}

\texttt{which\_better} generalizuje bolje nego \texttt{which\_algo} u oba scenarija (F1
$0.538$--$0.610$ naspram $0.360$--$0.481$) i ostaje blizu svoje 5-fold reference ($0.559$).
\texttt{which\_algo} pada ispod svoje reference ($0.402 -> 0.360$) baš na
interpolisanom \texttt{bf}$=7$, što je neobično: tačnost (0.487) ostaje razumna, ali F1
pada, jer model na tom specifičnom \texttt{bf} sistematski favorizuje većinsku klasu
Both na štetu ređih AB/MCTS predikcija. Kod ekstrapolacije na
\texttt{bf}$=14$ (ivica opsega), \texttt{which\_algo} F1 zapravo raste na $0.481$. Na
najvećim stablima je klasa Neither dominantnija i lakše prepoznatljiva, što
delimično kompenzuje otežanu ekstrapolaciju. Zajednička poruka: model ne uči trik vezan za
konkretnu vrednost \texttt{branching\_factor} (rezultati na neviđenom \texttt{bf} ne
odudaraju previše), ali apsolutni nivo tačnosti zavisi od toga koja mešavina klasa dominira na tom
delu opsega.

\subsection{Šta čini stablo teškim}

Klasa Neither (nijedan algoritam nije tačan) posebno se analizira binarnom
klasifikacijom Neither naspram ostalih (XGBoost, skup B), sa SHAP analizom važnosti.

\begin{figure}[H]
\centering
\includegraphics[width=0.95\textwidth]{../plots/visualizations/06_hard_tree_shap.png}
\caption{SHAP: prediktori teškog stabla (Neither naspram ostalih) i doprinos po dubini}
\end{figure}

\begin{table}[H]
\centering
\caption{Top 10 prediktora teškog stabla (SHAP, skup B)}
\begin{tabular}{lr}
\toprule
\textbf{Atribut} & \textbf{mean(|SHAP|)} \\
\midrule
branch\_top2\_gap       & 0.541 \\
ab\_reach\_depth        & 0.446 \\
branch\_mean\_spread    & 0.385 \\
branch\_value\_skew     & 0.312 \\
chance\_frac\_top1      & 0.285 \\
chance\_node\_count\_d1 & 0.203 \\
node\_budget            & 0.175 \\
top2\_size\_ratio       & 0.174 \\
budget\_coverage        & 0.173 \\
heur\_incoherence       & 0.169 \\
\bottomrule
\end{tabular}
\end{table}

\texttt{branch\_top2\_gap} je najvažniji, ali tesno praćen sa \texttt{ab\_reach\_depth}
(0.54 naspram 0.45): strukturna margina korena i realno dostupan budžet su, dakle,
podjednako jaki prediktori teškog stabla. Teška stabla su pre svega ona gde su prva dva
poteza iz korena vrednosno gotovo izjednačena i/ili gde AB nema dovoljno budžeta da dosegne
razumnu dubinu. Ostatak top deset atributa (\texttt{branch\_mean\_spread},
\texttt{branch\_value\_skew}, \texttt{chance\_frac\_top1}, \texttt{top2\_size\_ratio}) dodatno
opisuje oblik i asimetriju grana iz korena, dok \texttt{node\_budget},
\texttt{budget\_coverage} i \texttt{heur\_incoherence} potvrđuju da su budžet i kvalitet
heuristike nezavisni dobri prediktori.

%-------------------------------------------------------------------
\section{Zaključak}
%-------------------------------------------------------------------

\subsection{Sažetak nalaza}

\textbf{Nalaz 1.} Budžet dominira nad sirovom veličinom stabla, pritom modeli se ne slažu da li
ga koriste u sirovom ili normalizovanom obliku. Stablo odlučivanja deli koren
direktno po \texttt{node\_budget}, a XGBoost mu dodeljuje ubedljivo najveću
važnost. Slučajne šume, s druge strane, i dalje najviše teže strukturnoj margini
\texttt{branch\_top2\_gap}, sa logaritamski izvedenim \texttt{ab\_reach\_depth} (koliko nivoa
dubine budžet realno pokriva s obzirom na faktor grananja) tek na trećem mestu. Oba modela se
slažu da je budžet centralni signal, samo se razlikuju u tome kroz koji oblik ga najbolje
koriste.

\textbf{Nalaz 2.} Margina korena je nezavisan, podjednako jak izvor informacija.
\texttt{branch\_top2\_gap} (razlika između prva dva poteza iz korena, računata bez pokretanja
ijednog algoritma) je najvažniji pojedinačni prediktor kod Slučajnih šuma za
\texttt{which\_better} i ubedljivo najvažniji u SHAP analizi za to da li je stablo uopšte
teško rešivo (odeljak 6.4). 

\textbf{Nalaz 3.} \texttt{which\_better} je strukturno čistiji cilj od \texttt{which\_algo}.
Na istom skupu strukturnih atributa (B), \texttt{which\_better} dosledno nadmašuje
\texttt{which\_algo} (0.554 naspram 0.521 macro-F1, odeljak 5.6), bolje generalizuje na
neviđene vrednosti faktora grananja (odeljak 6.3) i ima uravnoteženiju raspodelu klasa
(odeljak 2.4). Razlog je strukturan, \texttt{which\_algo} meša koliko
je algoritam dobar sa koliko je stablo teško uopšte (klasa Both dominira kad god je
bilo koji uslov ispunjen), dok \texttt{which\_better} to razdvaja.

\textbf{Nalaz 4.} Struktura stabla je dovoljna za merljiv, ali ne potpun skup informacija. Skup B (28 atributa, isključivo strukturni) postiže macro-F1
$= 0.554$ za \texttt{which\_better}, daleko iznad slučajnog pogađanja ($\approx 0.33$) ali
nešto ispod skupa A koji uključuje merenja ponašanja algoritama (0.679). Ovi atributi
(posebno \texttt{mcts\_vote\_margin}) nose
stvarnu dodatnu vrednost, ali ne menjaju koji tip strukture je bitan.

\subsection{Ograničenja}

Stabla su sintetička i kontrolisanog oblika, realna stabla igara imaju bogatiju
strukturu i specifične evaluacione funkcije, pa nalazi na stvarnim
igrama nisu eksperimentalno proverena. Iscrpni minimax kao izvor istine nije primenjiv na igre
sa milijardama čvorova, pa bi na takvim igrama selekcija algoritma morala da se osloni na
aproksimativan algoritam. Međutim, svrha ovog projekta i jeste bila kontrolisanje strukture stabla
radi raznovrsnosti skupa podataka, i kasnija validacija na neviđenim stablima.

\end{document}
